\chapter{직접 돌려 보기}\label{app:run}

이 문서는 코드를 읽기 위한 것이지만, 읽은 것을 직접 돌려 보면 이해가 빠르다. 여기서는 어디서 무엇을 돌리는지만
짧게 적는다. 설치의 세부는 저장소의 \code{README.md}와 각 폴더의 \code{README.md}에 있다.

\begin{center}
\small
\begin{tabularx}{\linewidth}{>{\raggedright\arraybackslash}p{46mm}>{\raggedright\arraybackslash}X}
\toprule
하고 싶은 것 & 방법 \\
\midrule
봄베의 전 과정을 따라 하기 & 노트북 \code{notebooks/01\_crib\_to\_plaintext.ipynb}. 크립 정렬부터 체킹까지를
  단계별로 실행한다. 2부를 읽으며 함께 보기에 알맞다 \\
Rust 코어의 시험 & 저장소 루트에서 \code{cargo test} \\
전체 탐색의 성능 측정 & \code{cargo run -{}-release -p christopher-core -{}-example full\_search enigma-i} \\
거짓 정지의 통계 & \code{cargo run -{}-release -p christopher-core -{}-example stop\_statistics 400} \\
기준 RTL의 시험 & \code{rtl/tb/run.sh}. Verilator와 cocotb가 필요하다(\code{rtl/tb/README.md}).
  \code{-k scrambler}처럼 일부만 고를 수 있고, \code{WAVES=1}로 파형을 남길 수 있다 \\
하드웨어 보드의 생성과 검증 & 보드 폴더(예: \code{hardware/relay/shifter-board})에서 \code{python gen\_shifter.py},
  \code{python verify\_shifter.py}. KiCad 10의 \code{kicad-cli}가 필요하다. 보드마다의 순서는 그 폴더의 \code{README.md} \\
이 문서의 빌드 & \code{report/guide} 폴더에서 \code{latexmk} \\
\bottomrule
\end{tabularx}
\end{center}

3부를 읽을 때에는 파형을 남겨 놓고 \code{bombe\_core}의 \code{live\_q}가 클럭마다 넓어지는 것을 직접 보기를
권한다. 그림~\ref{fig:g-core-timing}이 실제 신호로 보인다.
